% !TeX root = ../../Book3.tex
% 纲要标题：### 15.2 Koszul 复形
% 纲要位置：工作记录/计划纲要.md，第 2455 行。
% 纲要细目（写作范围）：
% * 外代数上的 Koszul 微分
% * 一元情形与张量积构造
% * 符号约定、微分平方为零与乘法同伦；同调的局部化和支集、含单位理想时的退化
% * 生成元可逆线性变换诱导的 Koszul 复形同构

\section{Koszul 复形}
\label{b3:sec:1502}

一条乘法映射记录一个方程的核与商，多个方程则还需要记录它们之间的关系。外代数的符号使这些关系组成复形，其同调能同时反映不同阶的失效。


% 纲要内容（仅作写作计划，尚非教材正文）：
%
% * 外代数上的 Koszul 微分
% * 一元情形与张量积构造
% * 符号约定及微分平方为零
% * 生成元可逆线性变换诱导的 Koszul 复形同构
%

\subsection{外代数上的 Koszul 微分}
\label{b3:subsec:1502:01}

\begin{definition}\label{b3:def:koszul-complex}
设 \(R\) 为交换环，\(r\) 为非负整数，\(x=(x_1,\ldots,x_r)\in R^r\)，令 \(F=R^r\) 的标准基为 \(e_1,\ldots,e_r\)。取
\[
K_i(x;R)=\bigwedge^i_RF\quad(0\le i\le r),\qquad K_i=0\quad(i<0\;\text{或}\;i>r),
\]
并以 \(R\) 线性延拓规定
\begin{equation}\label{b3:eq:koszul-differential}
\mathrm{d}_i(e_{j_1}\wedge\cdots\wedge e_{j_i})
=\sum_{a=1}^i(-1)^{a-1}x_{j_a}
e_{j_1}\wedge\cdots\wedge\widehat e_{j_a}\wedge\cdots\wedge e_{j_i}.
\end{equation}
这里 \(j_1<\cdots<j_i\)，帽号表示删去该因子，\(\mathrm{d}_0=0\)。所得链复形记为 \(K_\bullet(x;R)\)，称为 \(x\) 的\term[koszul-fuxing]{Koszul 复形}[Koszul complex]。对 \(R\) 模 \(M\)，定义 \(K_\bullet(x;M)=K_\bullet(x;R)\otimes_RM\)。
\end{definition}
\symidx{\(K_\bullet(x;M)\)：Koszul 复形}

设 \(R\) 为交换环，\(C_\bullet\) 为 \(R\) 模的链复形。本章采用下标递降的微分 \(\mathrm{d}_i:C_i\rightarrow C_{i-1}\)，它满足 \(\mathrm{d}_{i-1}\mathrm{d}_i=0\)。核 \(Z_i=\ker \mathrm{d}_i\) 中的元素称为\term[quan]{圈}[cycle]，像 \(B_i=\operatorname{im}\mathrm{d}_{i+1}\) 中的元素称为\term[bianjie]{边界}[boundary]。微分的复合为零保证 \(B_i\subseteq Z_i\)，其商 \(H_i(C)=Z_i/B_i\) 为\term[tongdiao-mo]{同调模}[homology module]；它记录核中还有哪些元素不能由前一级的像解释。这与第四卷的上链记号 \(C^{-i}=C_i\) 相容；上链微分仍升次。本节先把所需链复形直接写出，微分平方为零将在下文核验。

对任意 \(R\) 模 \(M\)，Koszul 复形的次数一到次数零的映射为
\[
M^r\xrightarrow{(x_1\ \cdots\ x_r)}M.
\]
它的像是 \(IM\)，其中 \(I=(x_1,\ldots,x_r)\)，而次数零的微分为零，故总有
\[
H_0(x;M)=M/IM.
\]
这个零次同调就是施加全部方程后所得的商；正次数同调则记录 Koszul 复形在相应次数未能正合的部分。

\ProseParagraphBreak
比较两个复形时，各次数的映射还须尊重微分，才能把关于核与像的比较传到同调。具体地，设 \(R\) 为交换环，\(C_\bullet,D_\bullet\) 为 \(R\) 模的链复形。若一族 \(R\) 线性映射 \(f_i:C_i\to D_i\) 满足 \(\mathrm{d}_Df_i=f_{i-1}\mathrm{d}_C\)，则称 \(f=(f_i)\) 为从 \(C_\bullet\) 到 \(D_\bullet\) 的\term[lian-yingshe]{链映射}[chain map]。对圈 \(z\)，相容式给出 \(\mathrm{d}_Df_i(z)=0\)；对边界 \(\mathrm{d}_Cw\)，它给出 \(f_i(\mathrm{d}_Cw)=\mathrm{d}_Df_{i+1}(w)\)。所以圈映到圈、边界映到边界，从而诱导同调模之间的同态。

\subsection{一元情形与张量积构造}
\label{b3:subsec:1502:02}

一元情形为
\[
K_\bullet(x;M):\qquad 0\longrightarrow M\xrightarrow{x}M\longrightarrow0,
\]
左边的 \(M\) 位于次数一，右边位于次数零。因此
\[
H_1(x;M)=0:_Mx,\qquad H_0(x;M)=M/xM.
\]
二元情形的全部矩阵则是
\[
0\longrightarrow M\xrightarrow{\binom{-y}{x}}M^2
\xrightarrow{(x\ \ y)}M\longrightarrow0.
\]
首个矩阵来自 \(\mathrm{d}(e_1\wedge e_2)=xe_2-ye_1\)。两个矩阵的复合把 \(c\in M\) 送到 \(x(-yc)+y(xc)=0\)，所以二元情形的微分平方为零正是交换律的结果。次数一的圈是满足 \(xa+yb=0\) 的关系 \((a,b)\)，边界则是 \((-yc,xc)\) 这种由交换律直接产生的 Koszul 关系。因此
\[
H_1(x,y;M)
=\frac{\{(a,b)\in M^2:xa+yb=0\}}
       {\{(-yc,xc):c\in M\}},
\]
它保留不能由这些 Koszul 关系得到的其余关系。

把每个方程的两项复形一起张量，就得到多元 Koszul 复形。下面的符号约定保证张量微分与外积微分相容。

\begin{proposition}\label{b3:prop:koszul-tensor}
设 \(R\) 为交换环，\(r\) 为非负整数，\(x_1,\ldots,x_r\in R\)。约定任意两个 \(R\) 模链复形 \(C_\bullet,D_\bullet\) 的张量积在次数 \(n\) 取 \(\bigoplus_{p+q=n}C_p\otimes_RD_q\)，并对齐次元 \(u\in C_p,v\in D_q\) 规定
\[
\mathrm{d}(u\otimes v)=\mathrm{d}_Cu\otimes v+(-1)^p u\otimes \mathrm{d}_Dv,
\]
在此约定下有链复形同构
\[
K_\bullet(x_1,\ldots,x_r;R)
\cong K_\bullet(x_1;R)\otimes_R\cdots\otimes_RK_\bullet(x_r;R).
\]
\end{proposition}
\begin{proof}
上述张量微分连续作用两次时，\(\mathrm{d}_C^2\) 与 \(\mathrm{d}_D^2\) 所给的项为零，两个交叉项的符号分别为 \((-1)^{p-1}\) 与 \((-1)^p\)，相加也为零，因而它确实定义链复形。在二元情形中，\(e_1\otimes e_2\) 的第一个因子次数为一，所以
\[
\mathrm{d}(e_1\otimes e_2)
=x_1(1\otimes e_2)-x_2(e_1\otimes1),
\]
恰好对应 \(\mathrm{d}(e_1\wedge e_2)=x_1e_2-x_2e_1\)。一般地，在右边每个因子中选择次数零的 \(1\) 或次数一的 \(e_j\)。恰有 \(i\) 个次数一因子的张量给出次数 \(i\) 的一组基，按指标顺序送到相应外积。在第 \(a\) 个次数一因子上作用微分时，前面有 \(a-1\) 个次数一因子，张量微分的符号为 \((-1)^{a-1}\)，与\eqref{b3:eq:koszul-differential}一致。这给出各次数的双射并保持微分；\(r=0\) 时两边均为集中在次数零的 \(R\)。
\end{proof}

\subsection{符号约定与微分平方为零}
\label{b3:subsec:1502:03}

\begin{proposition}\label{b3:prop:koszul-signs-homotopy}
设 \(R\) 为交换环，\(r\) 为非负整数，\(x=(x_1,\ldots,x_r)\in R^r\)，\(M\) 为 \(R\) 模，\(I=(x_1,\ldots,x_r)\)。Koszul 微分满足 \(\mathrm{d}^2=0\)，并且 \(I\) 零化每个 \(H_i(x;M)\)。对每个整数 \(i\) 和每个含一的乘法集 \(S\subseteq R\)，有自然同构
\[
S^{-1}H_i(x;M)\cong H_i(x/1;S^{-1}M),
\]
其中右端在 \(S^{-1}R\) 上计算，\(x/1=(x_1/1,\ldots,x_r/1)\)。此外，
\[
\operatorname{Supp}_R H_i(x;M)
\subseteq\operatorname{Supp}_R M\cap V(I).
\]
若 \(I=R\)，则所有 \(H_i(x;M)\) 都为零；当 \(M\ne0\) 时，这样的序列不满足本书正则序列的约定。
\end{proposition}
\begin{proof}
二元计算中的两种次序在一般情形中也逐对出现。在一次外积上连续删去第 \(a\) 与第 \(b\) 个因子，其中 \(a<b\)，两种删法分别带符号
\[
(-1)^{a-1+b-2},\qquad (-1)^{b-1+a-1}.
\]
两项的系数同为 \(x_{j_a}x_{j_b}\)，符号相反，所以逐对相消。此论证在特征二也成立，因为此时两项之和为零。

令 \(h_j(w)=e_j\wedge w\)，并张量上 \(M\) 的恒等映射，得到次数升一的映射 \(h_j\)。外积微分满足 \(\mathrm{d}(e_j\wedge w)=x_jw-e_j\wedge\mathrm{d}w\)，移项便得
\[
\mathrm{d} h_j+h_j\mathrm{d}=x_j\operatorname{id}.
\]
若 \(\mathrm{d}w=0\)，便有 \(x_jw=\mathrm{d}(h_jw)\)，即 \(x_j\) 乘一个圈得到边界。因此 \(x_jH_i=0\)，任意 \(I\) 中元素也零化同调。

在外幂的标准基下，每个 \(K_i(x;M)\) 是有限个 \(M\) 的直和。这些外积基向量及微分公式经局部化，分别给出 \(K_i(x/1;S^{-1}M)\) 及其微分。由\cref{b3:thm:localization-exact}，局部化与微分的核、像和商相容，所以圈模与边界模局部化后仍是相应的圈模与边界模，再取商就得到所述同调同构。若 \(M_{\mathfrak p}=0\)，局部化后的复形各项都为零，故同调也为零。若 \(\mathfrak p\notin V(I)\)，某个 \(x_j/1\) 在 \(R_{\mathfrak p}\) 中可逆，又零化局部化后的同调，该同调只能为零。这证明支集包含于 \(\operatorname{Supp}_R M\cap V(I)\)，不需要有限生成假设。

若 \(1=\sum_j a_jx_j\)，取 \(h=\sum_j a_jh_j\)，则 \(\mathrm{d}h+h\mathrm{d}=\operatorname{id}\)。对任意圈 \(w\)，这个等式给出 \(w=\mathrm{d}(hw)\)，所以全部同调为零，包括 \(H_0=M/IM\)。但末端商为零，故在 \(M\ne0\) 时不能据此得到正则序列。例如，对域 \(k\)，在 \(A=k[X,Y]_{(X)}\) 中，\(Y\) 因为不属于 \((X)\) 而成为单位。因此 \(K_\bullet(X,Y;A)\) 的全部同调消失，而非零环 \(A\) 的末端商 \(A/(X,Y)A\) 为零，序列 \(X,Y\) 不正则。
\end{proof}

设 \(C_\bullet,D_\bullet\) 为交换环 \(R\) 上的链复形，\(f,g:C_\bullet\rightarrow D_\bullet\) 为链映射。若存在 \(R\) 线性映射 \(h_i:C_i\rightarrow D_{i+1}\)，使 \(f-g=\mathrm{d}_Dh+h\mathrm{d}_C\)，则称 \(f\) 与 \(g\) \term[lian-tonglun]{链同伦}[chain homotopy]。上面构造的 \(h_j\) 表明乘法 \(x_j\) 与零映射链同伦。链同伦映射在同调上相同：对圈 \(z\)，有 \((f-g)z=\mathrm{d}_D(hz)\)。

\ProseParagraphBreak
生成元的重排是可逆线性变换的一种特例。添加另一个生成元的倍数或用单位倍替换一个生成元，也有同样的复形不变性。

\begin{proposition}[生成元的可逆线性变换]\label{b3:prop:koszul-generator-change}
设 \(R\) 为交换环，\(M\) 为 \(R\) 模，\(r\) 为非负整数。把 \(x=(x_1,\ldots,x_r)^{\mathsf T}\) 看成列向量，取 \(U=(u_{ji})\in\operatorname{GL}_r(R)\)，令 \(y=Ux\)，即 \(y_j=\sum_i u_{ji}x_i\)。则有链复形同构
\[
K_\bullet(y;M)\xrightarrow{\sim}K_\bullet(x;M).
\]
特别地，置换生成元的次序不改变 Koszul 同调。
\end{proposition}
\begin{proof}
设 \(e_1,\ldots,e_r\) 是定义 \(K_\bullet(x;R)\) 所用的基，\(e'_1,\ldots,e'_r\) 是定义 \(K_\bullet(y;R)\) 所用的基。规定
\[
\varphi(e'_j)=\sum_i u_{ji}e_i.
\]
这个线性映射的矩阵为 \(U^{\mathsf T}\)，故可逆，而且
\[
\mathrm{d}_x\varphi(e'_j)=\sum_i u_{ji}x_i=y_j
=\mathrm{d}_y(e'_j).
\]
两边的微分都按外积的带符号 Leibniz 公式延拓，上式因此给出
\(\mathrm{d}_x\bigwedge\varphi=(\bigwedge\varphi)\mathrm{d}_y\)。各次外幂均可逆，再张量上 \(M\) 的恒等映射，即得到所述链复形同构。
\end{proof}

这个同构说明可逆线性变换保持 Koszul 复形及其同调，却不能在一般环与模上直接推出正则性也保持。前节的直和反例中，\(X,Y\) 与 \(Y,X\) 的 Koszul 复形同构，但只有第一个次序正则；正则性还须逐次检查商模上的乘法单射。
